% AI-written, AI cross-reviewed; no human sign-off; not a claim that RH is proved
% Atlas unit P014. Local manuscript v3, item 6.14, heading on printed PDF p. 58.
% UpperTheory.Top.every_character_of_green_v_win
% paper2/upper.tex:107–113. Verbatim excerpt; SHA-256 33d3df656d56d4b76bd5e32e7c60ce8d757a6e53f949df7bf8903d043d9cd7bd.
% Original TeX uses the macros and cross-reference keys of the source manuscript.


% Start of atlas edition annotation. The following block is outside the original excerpt.
\paragraph{Status of named inputs in this reading.}
\noindent\nolinkurl{GlobalTheory.GLq.Green} --- hypothesis at this head (\path{Development/AbelianVarietyTheory/Definitions/NonGaloisDescent/LabelDatum.lean:108--112}); variant.\par
% End of atlas edition annotation.

\begin{theorem}[Arbitrary characters from the three inputs]\label{pn:P014}
Assume $\mathcal B,V,\mathcal F$. Then $H$ holds, with the principal-character exception in P003.
\end{theorem}
\nodestatus{P014}
\begin{proof}
P015 supplies zero-freeness of the inducing primitive character. If $\chi$ is nonprincipal its inducing character cannot be the modulus-one character. Thus the change-of-level identity applies with the stated pole exception. Every factor in its finite Euler product is nonzero since $|\chi^*(p)p^{-s}|<1$. This gives the conclusion (blueprint Chapter~E, \emph{From primitive to arbitrary characters}~\bploc{prop:E-imprimitive}, \nolinkurl{CharacterBoxes/Imprimitive.lean:5}).
\end{proof}
